%% =============================================================================
%% main.tex — Content Entry Point (Completely Decoupled from style/healthcard.sty)
%% Compile with standard pdfLaTeX (NO -shell-escape required):
%%   latexmk -pdf main.tex      (outputs to build/main.pdf)
%%   make                       (or make preview / make clean)
%% =============================================================================
\documentclass[sigconf, nonacm]{acmart}

%% All visual styling, palette, cards, ORCID badges, and microtype live in style/.
%% One-click theme & style options:
%%   theme=pastel|ocean|sage|lavender|mono  (default: pastel = #5BCEFA/#F5A9B8/#FFFFFF)
%%   orcid=brand|official                   (default: brand)
%%   ribbon=gradient|stripe|none            (default: gradient)
%%   accent=blue|pink                       (default: blue)
%%   headerstyle=layered|clean              (default: layered)
\usepackage[theme=pastel]{style/healthcard}

%% Optional custom hex skinning (uncomment to override any 2 hex colors):
%% \hcsettheme{primary=5BCEFA, accent=F5A9B8}

%% -----------------------------------------------------------------------------
%% Logo Slots (Page 1 Dual 18pt Logos + Inner-Page 10pt Monochrome Mini Logo)
%% For custom image files in figures/, simply replace with:
%%   \projectlogoimage[18pt]{aminfer_logo.pdf}
%%   \orglogoimage[18pt]{srinternet_logo.pdf}
%%   \headerlogoimage[10pt]{aminfer_mini.pdf}
%% -----------------------------------------------------------------------------
%% =============================================================================
%% figures/logos.tex — Pure-TikZ Vector Demo Logos (100% pdfLaTeX / arXiv Safe)
%% Replace in main.tex with \projectlogoimage[18pt]{...} / \orglogoimage[18pt]{...}
%% whenever you have custom PDF/PNG brand marks in figures/.
%% =============================================================================

%% Project Logo (18pt height, adapts automatically to active theme colors)
\newcommand{\demoprojectlogo}{%
  \begin{tikzpicture}[baseline=0.3ex]
    \fill[hcPrimarySoft, rounded corners=4pt] (0,0) rectangle (18pt,18pt);
    \draw[hcPrimaryBorder, line width=0.55pt, rounded corners=4pt] (0,0) rectangle (18pt,18pt);
    \draw[hcPrimaryDark, line width=1.0pt, line cap=round, line join=round]
      (2.4pt,9pt) -- (5.5pt,9pt) -- (7.3pt,13.8pt) -- (10.1pt,4.2pt) -- (12.1pt,9pt) -- (15.4pt,9pt);
    \fill[hcAccent] (15.4pt,9pt) circle (1.5pt);
    \node[anchor=west, font=\sffamily\bfseries\fontsize{8.2}{9}\selectfont, text=hcPrimaryDark, inner sep=0pt]
      at (22pt, 9pt) {AuroraCard};
  \end{tikzpicture}%
}

%% Organization / Studio Logo (18pt height, adapts automatically to active theme)
\newcommand{\demoorglogo}{%
  \begin{tikzpicture}[baseline=0.3ex]
    \fill[hcSurface, rounded corners=4pt] (0,0) rectangle (18pt,18pt);
    \draw[hcAccentBorder, line width=0.6pt, rounded corners=4pt] (0,0) rectangle (18pt,18pt);
    \fill[hcPrimary] (4.5pt,4.5pt) rectangle (9.5pt,13.5pt);
    \fill[hcAccent] (8.8pt,6.8pt) rectangle (13.8pt,13.5pt);
    \node[anchor=west, font=\sffamily\bfseries\fontsize{7.8}{9}\selectfont, text=hcMuted, inner sep=0pt]
      at (22pt, 9pt) {Pastel Lab};
  \end{tikzpicture}%
}

%% Inner-Page 10pt Monochrome Miniature Logo (Pages 2+ Top-Left Running Header)
\newcommand{\demoheaderlogo}{%
  \begin{tikzpicture}[baseline=0.15ex]
    \draw[hcMuted, line width=0.65pt, rounded corners=2pt] (0,0) rectangle (10pt,10pt);
    \draw[hcMuted, line width=0.75pt, line cap=round, line join=round]
      (1.5pt,5pt) -- (3.3pt,5pt) -- (4.3pt,7.6pt) -- (5.8pt,2.4pt) -- (6.8pt,5pt) -- (8.5pt,5pt);
  \end{tikzpicture}%
}

\projectlogo{\demoprojectlogo}
\orglogo{\demoorglogo}
\headerlogo{\demoheaderlogo}

%% -----------------------------------------------------------------------------
%% Paper Metadata
%% -----------------------------------------------------------------------------
\reportlabel{Technical Report}
\reportbadge{v1.0 \textperiodcentered\ Preprint}
\reportdate{November 21, 2026}

\title{AuroraCard: Calibrated Multimodal Vital-Sign Fusion for Longitudinal Clinical Decision Support}
\subtitle{A Modular Architecture for Uncertainty-Aware Physiological Monitoring and Interpretable Risk Stratification}
\hcshorttitle{AuroraCard: Calibrated Multimodal Vital-Sign Fusion}
\hcshortauthors{Lin, Takahashi, Vance, and Zhou}

%% Authors: \hcauthor[affiliation/note marks]{Full Name}{ORCID (or empty)}
\hcauthor[1,*]{Aria Lin}{0000-0002-1825-0097}
\hcauthor[1,2,*]{Sora Takahashi}{0000-0001-9432-6104}
\hcauthor[2]{Elena Vance}{0000-0003-2719-8451}
\hcauthor[1,\dag]{Mina Zhou}{0000-0002-6510-3392}

%% Affiliations
\hcaffiliation[1]{Pastel Systems Research Lab}
\hcaffiliation[2]{Department of Biomedical Informatics, Pacific Institute of Technology}

%% Author Notes & Correspondence
\hcauthornote{\textsuperscript{*}Equal contribution}
\hcauthornote{\textsuperscript{\dag}Corresponding author}
\hccorrespondence{\hcemail{aria.lin@pastelsys.org}, \hcemail{mina.zhou@pastelsys.org}}

%% Resource Metadata Pills (icons auto-matched by label)
\hcmetadata[Code]{\hcghlink{Ariakage/tech-report-template}}
\hcmetadata[Models]{\hchflink{pastel-lab/auroracard-v1}}
\hcmetadata[Dataset]{\hchflink{datasets/pastel-lab/vitalbench-26}}
\hcmetadata[Project]{\href{https://pastelsys.org/auroracard}{pastelsys.org/auroracard}}

%% Keywords (automatically rendered as pastel pill tags inside the Header Card)
\hckeywords{Clinical Decision Support, Multimodal Fusion, Uncertainty Calibration, Vital-Sign Modeling, Preprint Systems}

%% -----------------------------------------------------------------------------
%% Document Body
%% -----------------------------------------------------------------------------
\begin{document}

%% Place abstract BEFORE \maketitle so it renders inside the Header Card
\begin{abstract}
%% =============================================================================
%% sections/00_abstract.tex
%% =============================================================================
Longitudinal clinical monitoring routinely combines asynchronous high-frequency physiological waveforms with sparse laboratory panels and structured electronic health records. Existing multimodal neural architectures often achieve strong discriminative performance on retrospective cohorts yet suffer from severe probability miscalibration, silent degradation under sensor dropout, and brittle cross-hospital generalization. We introduce \textbf{AuroraCard}, a lightweight, calibration-first multimodal fusion framework built around gated continuous-time state-space encoders and conformal temperature routing. Across four multi-center intensive care and outpatient benchmarks totaling 184{,}200 patient-stays, AuroraCard improves macro-AUROC to \hlblue{96.4\%} (+3.8 points over strong Transformer baselines), reduces Expected Calibration Error by \hlpink{61.5\%} to \hlpink{0.015}, and runs at \hlblue{14.2\,ms} per patient window on a single workstation GPU. We further establish finite-sample coverage guarantees under structured modality missingness and release full model checkpoints, evaluation harnesses, and reproducible LaTeX reporting artifacts.

\end{abstract}

\maketitle

%% =============================================================================
%% sections/01_introduction.tex
%% =============================================================================
\section{Introduction}
\label{sec:intro}

Modern clinical surveillance systems ingest continuous bedside telemetry alongside intermittent laboratory assays, medication administration logs, and triage notes~\cite{johnson2023mimic, rajpurkar2022ai}. Despite rapid progress in clinical foundation models~\cite{moor2023foundation}, deploying predictive risk models in real-time hospital workflows remains constrained by three tightly coupled operational challenges:

\begin{itemize}
  \item \textbf{Asynchronous missingness.} Wearable and bedside sensors experience irregular sampling rates, motion artifacts, and transient lead disconnections that violate synchronized sequence assumptions~\cite{shukla2021multi}.
  \item \textbf{Confidence miscalibration.} High-capacity multimodal Transformers frequently assign overconfident risk scores during distribution shifts, undermining clinician trust and triggering alarm fatigue~\cite{guo2017calibration, ovadia2019can}.
  \item \textbf{Latency and interpretability budgets.} Edge clinical workstations require sub-20\,ms inference per patient card while exposing modality-level attribution and calibrated uncertainty intervals.
\end{itemize}

To address these requirements simultaneously, we present \textbf{AuroraCard}, a modular physiological fusion architecture that treats each patient encounter as a structured, uncertainty-aware state representation. Rather than concatenating imputed dense tensors into a single quadratic attention stack, AuroraCard encodes each stream via a continuous-time decay kernel and fuses available modalities through an entropy-gated cross-modal router.

\begin{hccard}[Core Contributions]
  \begin{enumerate}
    \item \textbf{Modality-aware continuous encoding:} We formulate a closed-form irregular-interval state update that eliminates naive forward-fill imputation artifacts.
    \item \textbf{Conformal temperature routing:} A lightweight calibration layer adjusts predictive entropy dynamically based on real-time sensor completeness.
    \item \textbf{Multi-center validation:} Across four public and federated cohorts, AuroraCard achieves state-of-the-art discrimination, calibration, and edge latency.
  \end{enumerate}
\end{hccard}

Figure~\ref{fig:metrics} summarizes the headline operational gains achieved by AuroraCard across our primary multi-center test split, illustrating simultaneous improvements in discriminative accuracy, probability calibration, missing-modality robustness, and inference throughput.

\begin{figure*}[t]
  \begin{hcmetricgrid}[4]
    \hcmetriccard{blue}{heart}{Macro AUROC}{96.4\%}{+3.8\% vs. Transformer} &
    \hcmetriccard{pink}{shield}{Calibration ECE}{0.015}{$-$61.5\% error reduction} &
    \hcmetriccard{blue}{bolt}{Window Latency}{14.2\,ms}{$-$44\% vs. Multimodal ViT} &
    \hcmetriccard{pink}{check}{Dropout Robustness}{94.1\%}{At 40\% sensor drop}
  \end{hcmetricgrid}
  \caption{Headline clinical monitoring metrics for AuroraCard on the multi-center VitalBench-26 evaluation suite ($N = 46{,}050$ held-out stays). All improvements are statistically significant ($p < 0.001$, paired bootstrap).}
  \Description{Four side-by-side health-card metric tiles showing Macro AUROC 96.4 percent, Calibration ECE 0.015, Window Latency 14.2 ms, and Dropout Robustness 94.1 percent.}
  \label{fig:metrics}
\end{figure*}

%% =============================================================================
%% sections/02_methodology.tex
%% =============================================================================
\section{Methodology}
\label{sec:method}

Let a longitudinal patient encounter over horizon $[0, T]$ consist of $M$ heterogeneous physiological modalities $\mathcal{X} = \{\mathcal{S}^{(1)}, \dots, \mathcal{S}^{(M)}\}$, where each stream $\mathcal{S}^{(m)} = \{(t_k^{(m)}, \mathbf{x}_k^{(m)}, \mathbf{m}_k^{(m)})\}_{k=1}^{L_m}$ contains irregular timestamps $t_k^{(m)}$, measurement vectors $\mathbf{x}_k^{(m)} \in \mathbb{R}^{d_m}$, and binary observation masks $\mathbf{m}_k^{(m)} \in \{0,1\}^{d_m}$.

\subsection{Continuous-Time Vital Encoding}
\label{sec:method:encoder}

Following continuous-time neural dynamics~\cite{gu2022efficiently, kidger2020neural}, we evolve each modality latent state $\mathbf{h}_k^{(m)} \in \mathbb{R}^{d}$ between consecutive timestamps $\Delta t_k = t_k^{(m)} - t_{k-1}^{(m)}$ using an exponential decay gate parameterized by positive learnable rates $\boldsymbol{\gamma}_m = \operatorname{softplus}(\mathbf{w}_m)$:
\begin{align}
  \tilde{\mathbf{h}}_{k}^{(m)} &= \exp\!\big(-\boldsymbol{\gamma}_m \Delta t_k\big) \odot \mathbf{h}_{k-1}^{(m)}, \nonumber \\
  \mathbf{h}_k^{(m)} &= \operatorname{GRUCell}\!\big(\mathbf{u}_k^{(m)}, \tilde{\mathbf{h}}_k^{(m)}\big),
  \label{eq:ct-decay}
\end{align}
where $\mathbf{u}_k^{(m)} = \phi_m(\mathbf{x}_k^{(m)} \odot \mathbf{m}_k^{(m)}, \log(1 + \Delta t_k))$ projects observed channels and inter-arrival intervals into the shared latent space. When a channel remains unobserved for an extended interval, its latent influence decays smoothly toward the prior baseline without injecting artificial discontinuities.

\begin{hcdefinition}{Modality Reliability Score}{rel}
  Given the instantaneous observation density $\rho_m(t) \in [0, 1]$ over lookback window $[t - W, t]$ and encoder epistemic variance $\sigma_m^2(t)$, the \emph{reliability weight} of modality $m$ is defined as
  \begin{equation}
    \alpha_m(t) = \frac{\rho_m(t)\,\exp\!\big(-\beta \sigma_m^2(t)\big)}{\sum_{j=1}^{M} \rho_j(t)\,\exp\!\big(-\beta \sigma_j^2(t)\big) + \varepsilon},
    \label{eq:reliability}
  \end{equation}
  with temperature sensitivity $\beta > 0$ and numerical stabilizer $\varepsilon = 10^{-6}$.
\end{hcdefinition}

\subsection{Entropy-Gated Fusion and Calibration}
\label{sec:method:fusion}

Rather than using static modality weights, AuroraCard aggregates the modality embeddings $\mathbf{z}^{(m)} = \mathbf{h}_{L_m}^{(m)}$ into a unified patient-card representation $\mathbf{z}_{\mathrm{card}} = \sum_{m=1}^{M} \alpha_m(T)\,\mathbf{W}_m \mathbf{z}^{(m)}$ and modulates the output logit temperature $\tau(\mathcal{X})$ using the effective Shannon entropy $H(\boldsymbol{\alpha}) = -\sum_{m} \alpha_m \log \alpha_m$:
\begin{equation}
  \hat{p}(y \mid \mathcal{X}) = \operatorname{softmax}\!\left(\frac{\mathbf{W}_{\mathrm{cls}}\,\mathbf{z}_{\mathrm{card}}}{\tau_0 + \lambda \, \big(1 - \bar{\rho}(\mathcal{X})\big) + \eta H(\boldsymbol{\alpha})}\right),
  \label{eq:temp-routing}
\end{equation}
where $\bar{\rho}(\mathcal{X}) = \frac{1}{M}\sum_{m=1}^{M} \rho_m(T)$ measures aggregate sensor availability, and $(\tau_0, \lambda, \eta) \in \mathbb{R}_{>0}^3$ are fitted strictly on a held-out calibration split using proper scoring rules.

\begin{hctheorem}{Finite-Sample Conformal Coverage}{coverage}
  Let $\mathcal{D}_{\mathrm{cal}} = \{(\mathcal{X}_i, y_i)\}_{i=1}^{n}$ be exchangeable with a test encounter $(\mathcal{X}_{\mathrm{test}}, y_{\mathrm{test}})$ under arbitrary covariate-dependent sensor dropout $\mathcal{M}(\mathcal{X})$. Define nonconformity scores $s_i = 1 - \hat{p}(y_i \mid \mathcal{X}_i)$ from Eq.~\eqref{eq:temp-routing} and let $\hat{q}_{1-\delta}$ be the $\lceil (n+1)(1-\delta)\rceil / n$ empirical quantile of $\{s_i\}_{i=1}^{n}$. Then the prediction set $\mathcal{C}(\mathcal{X}_{\mathrm{test}}) = \{y : 1 - \hat{p}(y \mid \mathcal{X}_{\mathrm{test}}) \le \hat{q}_{1-\delta}\}$ satisfies
  \[
    1 - \delta \;\le\; \mathbb{P}\!\big(y_{\mathrm{test}} \in \mathcal{C}(\mathcal{X}_{\mathrm{test}})\big) \;\le\; 1 - \delta + \frac{1}{n + 1}.
  \]
\end{hctheorem}

Algorithm~\ref{alg:calibration} details the online conformal calibration pass executed once per deployment site, while Listing~\ref{lst:router} shows the core PyTorch router with key reliability and temperature lines highlighted.

\begin{hcalgorithm}[Online Conformal Calibration Pass]{label=alg:calibration}
  \Require Calibration split $\mathcal{D}_{\mathrm{cal}}$, trained encoder $f_\theta$, miscoverage $\delta$
  \Ensure Temperature tuple $(\tau_0, \lambda, \eta)$ and conformal quantile $\hat{q}_{1-\delta}$
  \State Fit $(\tau_0, \lambda, \eta)$ on $\mathcal{D}_{\mathrm{cal}}$ via L-BFGS \Comment{frozen backbone}
  \State Initialize nonconformity score buffer $\mathcal{S} \gets \emptyset$
  \For{each encounter $(\mathcal{X}_i, y_i) \in \mathcal{D}_{\mathrm{cal}}$}
    \State Compute reliability weights $\boldsymbol{\alpha}_i$ and density $\bar{\rho}_i$ via Eq.~\eqref{eq:reliability}
    \State Evaluate calibrated probability $\hat{p}(y_i \mid \mathcal{X}_i)$ via Eq.~\eqref{eq:temp-routing}
    \State Append score $s_i \gets 1 - \hat{p}(y_i \mid \mathcal{X}_i)$ to $\mathcal{S}$
  \EndFor
  \State $\hat{q}_{1-\delta} \gets \operatorname{Quantile}\!\big(\mathcal{S},\, \lceil (n+1)(1-\delta) \rceil / n\big)$
  \State \Return $(\tau_0, \lambda, \eta),\, \hat{q}_{1-\delta}$
\end{hcalgorithm}

\begin{hccode}[AuroraCard Router (PyTorch)]{language=Python, hlblue={16,17}, hlpink={21,22,23}}{label=lst:router}
import torch
import torch.nn as nn
import torch.nn.functional as F

class ConformalRouter(nn.Module):
    """Entropy-gated calibration router."""
    def __init__(self, dim: int, classes: int):
        super().__init__()
        self.head = nn.Linear(dim, classes)
        self.tau0 = nn.Parameter(torch.tensor(1.05))
        self.lam = nn.Parameter(torch.tensor(0.35))
        self.eta = nn.Parameter(torch.tensor(0.18))

    def forward(self, emb, rho, var, beta=2.0):
        # Reliability weights alpha_m (Eq. 2)
        w = rho * torch.exp(-beta * var)
        alpha = w / (w.sum(dim=-1, keepdim=True) + 1e-6)
        z_card = torch.einsum("bm,bmd->bd", alpha, emb)

        # Adaptive temperature scaling (Eq. 3)
        ent = -(alpha * alpha.clamp_min(1e-8).log()).sum(-1)
        tau = self.tau0 + self.lam * (1 - rho.mean(-1))
        tau = tau + self.eta * ent
        logits = self.head(z_card) / tau.unsqueeze(-1)
        return F.softmax(logits, dim=-1), alpha
\end{hccode}

%% =============================================================================
%% sections/03_experiments.tex
%% =============================================================================
\section{Experimental Evaluation}
\label{sec:experiments}

\begin{table*}[t]
  \centering
  \caption{Multi-center benchmark results on VitalBench-26 ($N = 46{,}050$ held-out stays, mean $\pm$ std over 5 seeds). Lower is better for ECE ($\downarrow$) and Latency ($\downarrow$); higher is better for AUROC ($\uparrow$) and AUPRC ($\uparrow$). Best results are highlighted with delta badges relative to Multimodal Transformer.}
  \label{tab:main-results}
  \begin{hctable}
    \begin{tabularx}{\textwidth}{@{}X c c c c c c@{}}
      \hctoprule
      \hctablehead
      \hcth{Model Architecture} & \hcth{Params} & \hcth{Macro AUROC ($\uparrow$)} & \hcth{AUPRC ($\uparrow$)} & \hcth{ECE ($\downarrow$)} & \hcth{40\% Drop ($\uparrow$)} & \hcth{Latency ($\downarrow$)} \\
      \hcmidrule
      GRU-D~\cite{che2018recurrent} & 14.2\,M & 89.7 $\pm$ 0.3 & 54.1 $\pm$ 0.6 & 0.052 & 81.4 $\pm$ 0.5 & 19.6\,ms \\
      \hcaltrow
      mTAND~\cite{shukla2021multi} & 17.8\,M & 91.4 $\pm$ 0.2 & 58.3 $\pm$ 0.4 & 0.044 & 85.2 $\pm$ 0.4 & 22.8\,ms \\
      S4-Med~\cite{gu2022efficiently} & 16.5\,M & 92.1 $\pm$ 0.2 & 60.5 $\pm$ 0.5 & 0.041 & 87.9 $\pm$ 0.4 & 16.4\,ms \\
      \hcaltrow
      Multimodal Transformer & 19.1\,M & 92.6 $\pm$ 0.2 & 61.9 $\pm$ 0.4 & 0.039 & 86.3 $\pm$ 0.5 & 25.4\,ms \\
      \hcmidrule
      \hcbestrow
      \textbf{AuroraCard (Ours)} \hcpinkbadge{SOTA} & \textbf{17.4\,M} & \textbf{96.4}\gain{+3.8} & \textbf{69.8}\gain{+7.9} & \textbf{0.015}\drop{$-$61.5\%} & \textbf{94.1}\gain{+7.8} & \textbf{14.2\,ms}\speedup{1.8$\times$} \\
      \hcbottomrule
    \end{tabularx}
  \end{hctable}
\end{table*}

We evaluate AuroraCard on \textbf{VitalBench-26}, a unified multi-center clinical benchmark spanning four hospitals (184{,}200 encounters total; 70\% train, 5\% calibration, 25\% held-out test). All baselines are trained under identical preprocessing, fixed random seeds, and matched parameter budgets ($\approx 18\text{M}$ parameters) on a single NVIDIA L40S GPU.

\begin{hcnote}[Benchmark Protocol \& Measurement Note]
  Latency is measured end-to-end with batch size~1 across 1{,}000 warm-started streaming windows (p50 wall-clock time, including irregular-timestamp decay projection and conformal calibration head).
\end{hcnote}

\subsection{Main Multi-Center Benchmark}
\label{sec:exp:main}

Table~\ref{tab:main-results} compares AuroraCard against recurrent imputation baselines (GRU-D~\cite{che2018recurrent}), irregular-time attention models (mTAND~\cite{shukla2021multi}), structured state-space encoders (S4-Med~\cite{gu2022efficiently}), and a late-fusion Multimodal Transformer. AuroraCard consistently attains the highest discrimination (\hlblue{96.4\%} Macro AUROC) while cutting Expected Calibration Error (ECE) by \hlpink{61.5\%} across both ICU mortality and acute decompensation endpoints. Notably, our continuous-time encoder runs \hlblue{1.8$\times$} faster than the Multimodal Transformer, comfortably meeting real-time bedside display refresh budgets.

\subsection{Robustness Under Progressive Sensor Dropout}
\label{sec:exp:dropout}

Real-world clinical telemetry frequently experiences partial lead loss during patient transport or nursing interventions. To simulate realistic degradation, we randomly drop between 0\% and 60\% of continuous waveform channels at inference time without retraining. Figure~\ref{fig:dropout} plots Macro AUROC trajectories across escalating dropout severity. While standard cross-modal attention deteriorates rapidly beyond 30\% missingness, AuroraCard smoothly reweights surviving sensors via Eq.~\eqref{eq:reliability} and sustains $>91\%$ AUROC at 60\% disconnection.

%% =============================================================================
%% figures/fig_dropout.tex — Standalone Native PGFPlots Figure (Single-Source)
%% =============================================================================
\begin{figure}[t]
  \centering
  \begin{hcframecard}
    \begin{tikzpicture}
      \begin{axis}[
        hcaxis,
        xlabel={\hcbisep{缺失传感器通道比例}{Missing Sensor Channels} (\%)},
        ylabel={Macro AUROC (\%)},
        xmin=0, xmax=60,
        ymin=67, ymax=99,
        xtick={0,15,30,45,60},
        ytick={70,75,80,85,90,95},
        legend style={at={(0.03,0.04)}, anchor=south west}
      ]
        \addplot coordinates {(0,96.4) (15,95.9) (30,94.9) (45,93.5) (60,91.4)};
        \addlegendentry{AuroraCard (Ours)}

        \addplot coordinates {(0,92.6) (15,90.8) (30,88.1) (45,84.4) (60,79.2)};
        \addlegendentry{Multimodal Transformer}

        \addplot coordinates {(0,92.1) (15,90.9) (30,89.2) (45,86.7) (60,83.1)};
        \addlegendentry{S4-Med}

        \addplot coordinates {(0,89.7) (15,87.1) (30,83.6) (45,79.1) (60,74.2)};
        \addlegendentry{GRU-D}
      \end{axis}
    \end{tikzpicture}
  \end{hcframecard}
  \hcbicaption{%
    传感器渐进脱落条件下 AUROC 退化曲线。即使 60\% 生理波形流断连，AuroraCard 仍维持 $>91\%$ Macro AUROC。%
  }{%
    Graceful degradation under progressive channel dropout. AuroraCard maintains $>91\%$ AUROC even when 60\% of waveform streams are disconnected.%
  }
  \Description{Line chart comparing Macro AUROC from 0 to 60 percent missing channels for AuroraCard, Multimodal Transformer, S4-Med, and GRU-D.}
  \label{fig:dropout}
\end{figure}


\subsection{Component Ablation Analysis}
\label{sec:exp:ablation}

Table~\ref{tab:ablation} isolates the contribution of each architectural module. Removing continuous-time decay degrades missing-data discrimination by 2.9 points, while replacing conformal temperature routing with static temperature scaling nearly doubles calibration error on out-of-hospital shifts.

\begin{table}[t]
  \centering
  \caption{Ablation study on VitalBench-26 isolating continuous decay, reliability gating, and adaptive temperature routing.}
  \label{tab:ablation}
  \begin{hctable}
    \begin{tabularx}{\linewidth}{@{}X c c@{}}
      \hctoprule
      \hctablehead
      \hcth{Configuration Variant} & \hcth{AUROC ($\uparrow$)} & \hcth{ECE ($\downarrow$)} \\
      \hcmidrule
      \hcbestrow
      \textbf{AuroraCard (Full)} & \textbf{96.4} & \textbf{0.015} \\
      w/o Continuous Decay (Eq.~\ref{eq:ct-decay}) & 93.5 & 0.024 \\
      \hcaltrow
      w/o Reliability Gating (Eq.~\ref{eq:reliability}) & 94.2 & 0.029 \\
      w/o Adaptive Temperature (Eq.~\ref{eq:temp-routing}) & 96.3 & 0.031 \\
      \hcaltrow
      Static Mean-Imputed Baseline & 91.2 & 0.047 \\
      \hcbottomrule
    \end{tabularx}
  \end{hctable}
\end{table}

\begin{hcpinkcard}[Key Ablation Insight]
  Adaptive temperature routing leaves ranking performance virtually unchanged (96.3\% vs. 96.4\% AUROC) while reducing calibration error by \hlpink{51.6\%} (0.031 $\rightarrow$ 0.015). Decoupling uncertainty scaling from feature extraction thus eliminates the traditional accuracy--calibration trade-off.
\end{hcpinkcard}

\subsection{Cross-Hospital Conformal Coverage}
\label{sec:exp:coverage}

To verify Theorem~\ref{thm:coverage} empirically under site shift, we calibrate $\hat{q}_{1-\delta}$ at target miscoverage $\delta = 0.05$ on three source hospitals and evaluate coverage on the held-out fourth center~\cite{angelopoulos2023conformal, vovk2005algorithmic}. AuroraCard achieves \hlblue{95.3\%} empirical coverage with a mean prediction set size of 1.12 classes, compared to 88.4\% empirical coverage for uncalibrated softmax thresholding.

%% =============================================================================
%% sections/04_conclusion.tex
%% =============================================================================
\section{Conclusion and Future Work}
\label{sec:conclusion}

We introduced \textbf{AuroraCard}, a calibration-first multimodal fusion framework tailored to asynchronous physiological telemetry and longitudinal clinical surveillance. By combining continuous-time latent decay with reliability-weighted fusion and conformal temperature routing, AuroraCard delivers strong discriminative accuracy, verifiable finite-sample uncertainty guarantees, and sub-15\,ms edge inference latency.

\begin{hccaveat}[Clinical Deployment Caveat]
  While VitalBench-26 spans four geographically diverse medical centers, prospective bedside deployment still requires site-specific threshold governance, human-in-the-loop alert escalation policies, and continuous subgroup fairness auditing~\cite{rajpurkar2022ai}. All retrospective cohorts were de-identified under institutional ethics waivers.
\end{hccaveat}

\paragraph{Future Directions}
Promising extensions include incorporating structured clinician progress notes via retrieval-augmented adapters, federated conformal calibration across privacy-preserving hospital consortia, and streaming hardware quantization for wearable triage monitors.

\section*{Acknowledgments}
We thank the clinical informatics teams across participating partner hospitals for de-identified cohort curation, and the anonymous reviewers for constructive feedback on the conformal coverage analysis.

%% Automatically generated Author ORCID directory from \hcauthor declarations
\makeorcidlist

%% Reproducibility / Hardware Testbed Card (balances final column beside References)
%% Tip: For Apple Silicon / systems benchmarks (e.g. aminfer), swap rows to:
%%   SoC & Memory / OS & Kernel / Compiler & Backend (e.g. Zig & Metal) / Precision.
\begin{hcenvcard}[Hardware \& Evaluation Environment]
  \hcenvrow{Compute Node}{NVIDIA L40S (48\,GB) \textperiodcentered\ Edge: Apple M4}
  \hcenvrow{Host OS}{Ubuntu 24.04 LTS \textperiodcentered\ macOS 15.x}
  \hcenvrow{Toolchain}{PyTorch 2.6 / CUDA 12.6 \textperiodcentered\ Metal 3.2}
  \hcenvrow{Precision}{BF16 train \textperiodcentered\ FP16/FP32 calibrated eval}
  \hcenvrow{Reproducibility}{5 fixed seeds (42--46) \textperiodcentered\ Deterministic}
\end{hcenvcard}


%% Balance final page columns cleanly before References
\hcbalance

%% Bibliography (TeX Live built-in ACM-Reference-Format)
\bibliographystyle{ACM-Reference-Format}
\bibliography{bib/references}

\end{document}
